% ================================================================================================ capítulo 5
\chapter{Modo Imágenes}\label{cap:imagenes}

\begin{resumen}
\textbf{En este capítulo:} el \menu{Visor de píxeles} (escalas, regiones, cabecera FITS), la medición del FWHM y del
seeing de una estrella, y la calibración de una noche completa con \menu{Darks y flats}.
\end{resumen}

% ------------------------------------------------------------------------------------------------ visor
\section{Visor de píxeles}\label{sec:visor}

El visor abre imágenes FITS (también cubos de varios cuadros), PNG, TIFF de 16 bits, JPG y arreglos de NumPy, y funciona
como SAOImage DS9: cambias cómo se ve la imagen sin tocar sus datos, y lees el valor exacto de cada píxel. Puedes abrir
una imagen con \botonp{Abrir imagen…} o arrastrarla a la ventana.

\begin{capturam}{visor}{El visor con un campo de estrellas y una región rectangular.}
\item \textbf{Escala}: lineal, logarítmica, raíz cuadrada, asinh o histograma ecualizado.
\item \textbf{Límites}: mínimo–máximo, percentiles (99,5, 99 y 95\,\%), ZScale (el de DS9) o manual.
\item \textbf{Colores}: grises, grises invertidos, calor, viridis o frío–cálido.
\item \textbf{Forma de la región}: rectángulo, círculo o elipse.
\item \textbf{La imagen}. Rueda: zoom alrededor del cursor. Arrastrar: mover. \tecla{Mayús}\,+\,arrastrar o clic derecho
y arrastrar: dibujar una región. \boton{Ajustar} y \boton{1:1} cambian el zoom.
\item \textbf{Posición y valor} del píxel bajo el cursor (columna y fila desde 1; la fila 1 es la de arriba), y los
valores del negro y del blanco.
\item \textbf{Estadística de la región}: píxeles, media, mediana, σ, σ robusta, mínimo, máximo y suma, como en DS9.
\boton{Exportar CSV…} guarda la posición y el valor de cada píxel de la región.
\item \textbf{Cabecera FITS} completa.
\end{capturam}

La imagen se puede guardar en FITS de 32 bits o en PNG de 16 bits con \boton{Guardar FITS…} y \boton{Guardar PNG…}. En
un cubo, el número de cuadro aparece en la barra de arriba. El zoom no se aleja más allá de la imagen completa: siempre
la ves entera o más grande.

\subsection{Escala como en DS9}

La tarjeta \campo{Escala (como en DS9)} muestra el histograma de los valores de la imagen, en escala logarítmica, con dos
líneas: el \textbf{negro} (todo lo que esté por debajo se ve negro) y el \textbf{blanco}. Escribir otros valores en
\campo{Negro} y \campo{Blanco} cambia el contraste al instante y pone los límites en «Manual»; \boton{Automático} vuelve
a ZScale. Con \campo{Mismos límites en todos los cuadros}, al pasar de un cuadro a otro de un cubo los límites no cambian:
es la forma correcta de comparar dos imágenes, porque con ZScale cada una se ajusta por separado y una diferencia pequeña
desaparece.

\subsection{Medir una estrella: FWHM y seeing}\label{sec:fwhm_app}

Dibuja una región (mejor circular) alrededor de una estrella aislada y que no esté saturada, y pulsa
\botonp{Medir la estrella de la región}. TurbuLENS ajusta una gaussiana de dos dimensiones con fondo
(sección~\ref{sec:psf}) y dibuja sobre la imagen dos elipses: la verde a un sigma del centro, donde el brillo cae al
61\,\% del pico, y la amarilla a medio FWHM, donde cae a la mitad.

\begin{capturam}{visor_estrella}{Escala variable, medición de una estrella y su perfil radial.}
\item \textbf{Histograma} de la imagen con el negro (azul) y el blanco (rojo).
\item \campo{Negro} y \campo{Blanco}: los límites, a mano.
\item \campo{Mismos límites en todos los cuadros}.
\item \textbf{Escala de placa}, escrita o tomada con \boton{De Parámetros}.
\item \botonp{Medir la estrella de la región}.
\item \textbf{Resultado} del ajuste.
\item \textbf{Qué graficar}: perfil radial, histograma de la región, perfil horizontal o vertical.
\item \textbf{La gráfica}: aquí, el perfil radial con los píxeles, su promedio por anillos, la gaussiana ajustada y las
líneas de σ y de FWHM/2.
\end{capturam}

El resultado da el centro de la estrella con centésimas de píxel, el pico sobre el fondo, $\sigma_x$ y $\sigma_y$, la
elipticidad, el radio al primer sigma y el FWHM en píxeles. Con la escala de placa, el FWHM pasa a segundos de arco:
\[
\text{FWHM}\,('') = \text{FWHM}\,(\text{px}) \times s\,(''/\text{px}) .
\]

\begin{cuidado}[¿Cuándo el FWHM es el seeing?]
En una \textbf{exposición larga} (varios segundos) y \textbf{bien enfocada}, el tamaño de una estrella lo pone la
atmósfera y el FWHM es el seeing. En un cuadro de video de milisegundos, no: la estrella está «congelada» y sale más
pequeña. Tampoco si la imagen está desenfocada, si el seguimiento se corrió o si la estrella está saturada (la cima
plana engaña al ajuste). Por eso el DIMM es más confiable: mide movimientos, no tamaños. Las dos medidas juntas, en la
misma noche, son una buena comprobación.
\end{cuidado}

Si la región es muy pequeña, si solo tiene fondo o si la estrella es más grande que la región, el ajuste no se hace y el
mensaje dice qué cambiar.

% ------------------------------------------------------------------------------------------------ calibración
\section{Darks y flats}\label{sec:calibracion_app}

Una noche de observación deja muchas imágenes: las de los objetos (\textbf{luces}) y las de calibración (\textbf{bias},
\textbf{darks} y \textbf{flats}). Esta página las clasifica, decide qué calibración necesita cada luz y la aplica, todo a
partir de las cabeceras. Funciona con imágenes sueltas o con carpetas enteras, subcarpetas incluidas.

\begin{capturam}{calibracion_plan}{Una noche con luces de 30~s en filtro R y de 60~s en filtro V, darks de 30 y 120~s,
bias y flats de los dos filtros.}
\item \boton{Agregar carpeta…} y \boton{Agregar imágenes…}; también se pueden soltar en la ventana.
\item \textbf{La tabla}: cada archivo con su tipo, de dónde salió el tipo, exposición, ganancia o ISO, temperatura,
filtro, tamaño y binning.
\item \campo{márcala como}: corrige el tipo de las imágenes seleccionadas.
\item \textbf{Combinar con}: mediana (recomendado), media con recorte a 3~σ o media.
\item \campo{Escalar darks con el bias si no hay de la misma exposición}.
\item \textbf{El plan}: qué se resta y qué se divide en cada grupo de luces, con los avisos.
\item \botonp{Crear maestros y calibrar}.
\end{capturam}

\subsection{1. Imágenes}

TurbuLENS lee solo la \textbf{cabecera} de cada archivo (no la imagen entera, así que cientos de archivos se leen en
segundos): el tipo (\codigo{IMAGETYP}, \codigo{FRAME}), la exposición (\codigo{EXPTIME}, \codigo{EXPOSURE}), la ganancia o
el ISO, la temperatura (\codigo{CCD-TEMP}), el filtro, el binning, el tamaño y la cámara. En los TIFF y JPG de cámaras de
fotos lee los datos EXIF.

Si un archivo no dice qué es, el tipo sale del nombre del archivo («dark\_30s\_001») o de la carpeta («Darks/»,
«Flats/»). La columna \campo{Según} dice de dónde salió; lo que no se pudo saber se toma como luz y se marca en amarillo.
Si alguna quedó mal clasificada, selecciónala y cámbiala con \campo{márcala como}.

\subsection{2. Qué se va a hacer}

El plan agrupa las luces por exposición, filtro, ganancia y tamaño, y dice para cada grupo:
\begin{itemize}
\item \textbf{qué se resta}: el dark de la misma exposición; si no lo hay, un dark de otra exposición escalado con el
bias (sección~\ref{sec:dark_escalado}); si tampoco, solo el bias; o nada;
\item \textbf{qué se divide}: el flat del mismo filtro, al que antes se le restó su propio dark o el bias.
\end{itemize}
Los avisos salen en amarillo: darks a otra temperatura, luces sin exposición en la cabecera, falta de bias para
escalar\ldots{} En la captura, las luces de 60~s no tienen darks de 60~s y el plan escala el de 120~s por 0,5.

\subsection{3. Calibrar}

Elige la carpeta de salida y pulsa \botonp{Crear maestros y calibrar}. Cada maestro se crea una sola vez y se guarda en
\ruta{maestros/}; cada luz calibrada, en \ruta{calibradas/}, con su cabecera original más \codigo{CALSTAT},
\codigo{NDARK}, \codigo{NFLAT} y, si se escaló el dark, \codigo{DARKSCAL}. Las series grandes se combinan por franjas, con
un archivo temporal, sin llenar la memoria. Puedes cancelar en cualquier momento.

\captura{calibracion_resultado.png}{Al terminar: los maestros creados, el diagnóstico de cada flat y los botones para
revisar el resultado.}

Al terminar aparecen:
\begin{itemize}
\item \boton{Ver dark maestro}, \boton{Ver flat maestro}, \boton{Ver bias maestro} y \boton{Ver la primera calibrada},
que abren cada imagen en el visor.
\item \boton{Comparar antes y después}: abre la primera luz original y la calibrada como un cubo de dos cuadros, con
\campo{Mismos límites en todos los cuadros} activado. Pasa de un cuadro al otro y verás la corrección.
\item \boton{Usar en el modo medición}: pasa los maestros del grupo principal a la calibración de \menu{Parámetros}.
\end{itemize}

\subsection{El diagnóstico del flat}

Para cada flat, TurbuLENS informa cuánto corrige: «el flat normalizado va de 0,675 a 1,186; la corrección cambia el
brillo hasta un 32,5\,\%». Y avisa en dos casos:
\begin{itemize}
\item \textbf{Flats saturados.} Si muchos píxeles están en el valor máximo, el flat sale casi plano y no corrige nada.
Hay que tomarlos con menos exposición, con el histograma cerca de la mitad del rango.
\item \textbf{Flat casi uniforme.} Si varía menos del 1\,\%, la corrección es imperceptible. Es normal con un
telescopio sin viñeteo ni polvo, pero vale la pena revisarlo.
\end{itemize}

\begin{nota}[«El flat no cambia nada».]
Es la duda más común. Casi siempre tiene una de dos causas: los flats estaban saturados (el diagnóstico lo dice) o se
compararon las imágenes con ZScale, que ajusta el contraste de cada una por separado y esconde una corrección de pocos
por ciento. Con \boton{Comparar antes y después} los límites son los mismos y la diferencia se ve.
\end{nota}

\begin{consejo}[Solo calibraciones.]
Si agregas solo bias, darks o flats, sin luces, TurbuLENS crea los maestros (bias, darks por exposición y flats por
filtro) y los deja listos para otra noche.
\end{consejo}
